\documentclass{article}

\usepackage[utf8]{inputenc}     % pdfLaTeX
\usepackage[T1]{fontenc}
\usepackage[magyar]{babel}

\usepackage{amsthm}
\usepackage{amsmath}
\usepackage{amsfonts}
\usepackage{tikz}
\usetikzlibrary{arrows.meta, positioning}
\usepackage{graphicx}

\title{Solving combinatorial problems with code}
\author{Robin Rovenszky}
\date{Last updated: March 23, 2026}

\begin{document}

\maketitle
\tableofcontents
\newpage
\section{First problem}
In how many ways can we write ten as the sum of three non-negative integers, if the order of the addends matters?

\section{Second problem}
In how many ways can twenty be written as the sum of three positive integers, if the order of the addends matters?

In general, a number $n \in \mathbb{N}^+$ can be written as the sum of $k \in \mathbb{N}^+$ numbers in
\[
\binom{n-1}{k-1}
\]
ways, where $1 \le k \le n$.\\\\
If $k \in \mathbb{N}$, then it can be written in
\[
\binom{n + k - 1}{k - 1}
\]
ways.

\end{document}